% 数理逻辑（综述）
% license CCBYSA3
% type Wiki

本文根据 CC-BY-SA 协议转载翻译自维基百科\href{https://en.wikipedia.org/wiki/Mathematical_logic}{相关文章}。

数学逻辑（Mathematical Logic）是数学中对形式逻辑的研究。其主要分支包括模型论、证明论、集合论与递归论（又称可计算性理论）。数学逻辑的研究通常探讨逻辑形式系统的数学性质，例如其表达能力与推理能力。然而，它也可以用于刻画数学推理的正确性，或为数学建立理论基础。

自其诞生以来，数学逻辑既推动了数学基础研究的发展，也受到其驱动。该领域的起源可追溯至19世纪晚期，当时几何、算术与分析的公理化体系相继建立。20世纪初，大卫·希尔伯特（David Hilbert）提出了著名的“希尔伯特计划”（Hilbert’s Program），旨在证明数学基础理论的一致性。库尔特·哥德尔（Kurt Gödel）、格哈德·根岑（Gerhard Gentzen）等人的研究部分地解决了这一计划，同时也揭示了证明一致性问题的复杂性。

集合论的研究表明，几乎所有通常的数学内容都可以用集合语言形式化，尽管仍存在一些定理无法在常见集合论公理系统内被证明。现代数学基础研究的重点，往往不再是寻找能涵盖全部数学的理论，而是致力于分析哪些数学部分能够在特定的形式系统中被形式化（如“反向数学”所做的那样）。
\subsection{分支领域与研究范围}
1977年出版的《数学逻辑手册》（Handbook of Mathematical Logic由 J. Barwise 编辑，North-Holland 出版社，1977年\(^\text{[1]}\)）将当代数学逻辑大致划分为四个主要领域：
\begin{enumerate}
  \item 集合论（Set Theory）；
  \item 模型论（Model Theory）；
  \item 递归论（Recursion Theory）；
  \item 证明论与构造数学（Proof Theory and Constructive Mathematics），两者通常被视为同一研究方向的不同方面。
\end{enumerate}
此外，计算复杂性理论（Computational Complexity Theory）有时也被归入数学逻辑的研究范畴。\(^\text{[2][3]}\)
每个分支领域都有其独特的研究重点，但许多技术与结果在不同领域之间是共享的。各子领域之间的边界，以及数学逻辑与其他数学领域之间的界限，往往并不清晰。例如，哥德尔不完备定理不仅是递归论与证明论的重要成果，还引出了模态逻辑中的勒布定理（Löb's Theorem）；而强制法（Forcing）这一方法被广泛应用于集合论、模型论与递归论之中，并用于研究直觉主义数学。

范畴论（Category Theory）虽使用了大量形式公理化方法，并研究“范畴逻辑”（Categorical Logic），但通常不被视为数学逻辑的一个分支。由于范畴论在数学多个领域中具有广泛适用性，包括桑德斯·麦克莱恩（Saunders Mac Lane）在内的数学家提出将其作为独立于集合论的数学基础体系。这类基础体系依赖于“拓扑斯”（Topos）的概念，它可以被看作集合论的广义模型，并可在经典或非经典逻辑下定义。
\subsection{历史}
数学逻辑在19世纪中期作为数学的一个分支逐渐形成，源自哲学形式逻辑与数学传统的汇合。\(^\text{[4]}\)数学逻辑（又称“后勤学”（Logistic）、“符号逻辑”（Symbolic Logic）、“逻辑代数”（Algebra of Logic），或更近代的称呼“形式逻辑”（Formal Logic））是19世纪间发展起来的一系列逻辑理论的总称，它依赖一种人工符号系统与严格的演绎方法。\(^\text{[5]}\)在此之前，逻辑的研究主要与修辞学、计算术（calculationes\(^\text{[6]}\)）、三段论以及哲学研究相结合。20世纪上半叶，数学逻辑领域取得了爆炸性的发展，涌现出一系列奠基性成果，同时伴随着对数学基础问题的激烈争论。
\subsubsection{早期历史}
逻辑理论在历史上的许多文化中都得到了独立发展，包括古代中国、印度、希腊、罗马帝国以及伊斯兰世界。希腊的逻辑方法，尤其是亚里士多德的逻辑（即“项逻辑”，term logic），在其著作《工具论》（Organon）中系统呈现，对西方科学与数学的发展产生了长达数千年的深远影响。\(^\text{[7]}\)斯多葛派哲学家，尤其是克里西波斯（Chrysippus），开创了命题逻辑（propositional logic）的研究。18世纪的欧洲，一些哲学数学家——如莱布尼茨（Gottfried Wilhelm Leibniz）与兰伯特（Johann Heinrich Lambert）——尝试以符号化或代数化的方式处理形式逻辑的运算。然而，他们的研究成果大多零散，影响有限，并未引起广泛关注。
\subsubsection{十九世纪}
19世纪中叶，乔治·布尔（George Boole）与奥古斯都·德摩根（Augustus De Morgan）分别提出了系统化的数学逻辑理论。他们的研究基于代数学家乔治·皮科克（George Peacock）等人的工作，将传统的亚里士多德逻辑扩展为一个足以探讨数学基础的形式框架。\(^\text{[8]}\)1847年，瓦特罗斯拉夫·贝尔蒂奇（Vatroslav Bertić）独立于布尔，对逻辑的代数化也作出了重要贡献。\(^\text{[9]}\)此后，查尔斯·桑德斯·皮尔士（Charles Sanders Peirce）在布尔的成果基础上发展出一个包含关系与量词的逻辑体系，并于1870年至1885年间陆续发表了多篇论文。

哥特洛布·弗雷格（Gottlob Frege）则在1879年发表了《概念文字》（Begriffsschrift），独立提出了含量词的逻辑体系。这部著作通常被视为逻辑史上的分水岭。然而，弗雷格的工作在当时鲜为人知，直到20世纪初伯特兰·罗素（Bertrand Russell）开始推广其思想后才逐渐受到重视。弗雷格所创的二维逻辑符号体系虽然具有高度表达力，但并未被广泛采用，现代文献中几乎不再使用。

1890年至1905年间，恩斯特·施罗德（Ernst Schröder）出版了三卷本《逻辑代数讲义》（Vorlesungen über die Algebra der Logik）。  
该著作系统总结并扩展了布尔、德摩根与皮尔士的成果，是19世纪末符号逻辑研究的全面总结与重要参考文献。
